\documentclass{article}
\usepackage{graphicx}
\providecommand{\institute}[1]{}


\title{Provenance Decomposition of Pion Production in Pythia 8}
\subtitle{Latex Galore}

\author{CAP parton-tracking}
\institute{Wayne State University}
\date{\today}


\begin{document}


\maketitle
\begin{abstract}
This report presents a provenance decomposition of pion production in Pythia 8.  Every final-state hadron is traced through a generator-agnostic event-history graph back to its parent partons, separating the yield and the two-particle correlation into the parts that originate in the partonic stage, in hadronization, and in subsequent decays.
\end{abstract}




\section{Introduction}


Final-state hadron observables mix several physical origins: direct hadronization products, resonance feed-down, and weak-decay products.  This analysis attaches a full provenance record to every hadron so that each contribution can be measured separately.




\section{Method}


For every event the Pythia record is converted into an EventHistory directed graph.  A ProvenanceTagger walks each final hadron back up the graph, recording its production stage, its parent hadron (and whether that parent is a short-lived resonance), and its ancestor partons.  Single-particle spectra and two-particle correlations are then accumulated, decomposed by those provenance classes.




\section{Configuration and binning}


\begin{table}[htbp]
\centering
\begin{tabular}{l l}
\hline
parameter & value \\
\hline\hline
events & 20000 \\
$\sqrt{s}$ & 13000 GeV \\
QCD process & soft \\
random seed & 12345 \\
species (|pdg|) & 211 \\
\hline
\end{tabular}
\caption{Monte-Carlo run configuration.}
\label{tab:config}
\end{table}


\begin{table}[htbp]
\centering
\begin{tabular}{l r l}
\hline
observable & bins & range \\
\hline\hline
$p_{T}$ & 100 & 0 -- 10 GeV \\
$\eta$ & 120 & $-6$ -- $6$ \\
$\Delta\eta$ & 80 & $-8$ -- $8$ \\
$\Delta\phi$ & 72 & $-\pi$ -- $\pi$ \\
\hline
\end{tabular}
\caption{Analysis binning.}
\label{tab:binning}
\end{table}




\section{Single-particle results}


\begin{table}[htbp]
\centering
\begin{tabular}{l r r}
\hline
origin class & count & fraction \\
\hline\hline
Primary & 587869 & 31.67\% \\
FromResonance & 1071302 & 57.72\% \\
FromWeakDecay & 196763 & 10.60\% \\
Unknown & 0 & 0.00\% \\
\hline
\end{tabular}
\caption{Pion yield by production origin.}
\label{tab:origin}
\end{table}


\begin{table}[htbp]
\centering
\begin{tabular}{l r r}
\hline
parton class & count & fraction \\
\hline\hline
LightQuark & 1435838 & 77.36\% \\
Strange & 235015 & 12.66\% \\
Charm & 58315 & 3.14\% \\
Bottom & 4895 & 0.26\% \\
Gluon & 90202 & 4.86\% \\
NonPartonic & 31669 & 1.71\% \\
\hline
\end{tabular}
\caption{Pion yield by ancestor-parton flavour.}
\label{tab:parton}
\end{table}

\begin{figure}
\includegraphics[scale=0.75,height=2in]{}
\caption{Pion transverse-momentum spectrum, decomposed by production origin (primary hadronization product vs resonance / weak-decay feed-down).}
\label{}
\end{figure}
\begin{figure}
\includegraphics[scale=0.75,height=2in]{}
\caption{Pion pseudorapidity distribution, decomposed by production origin.}
\label{}
\end{figure}
\begin{figure}
\includegraphics[scale=0.75,height=2in]{}
\caption{Pion transverse-momentum spectrum grouped by the flavour of the ancestor pre-hadronization parton.}
\label{}
\end{figure}



\section{Two-particle results}


\begin{table}[htbp]
\centering
\begin{tabular}{l r r}
\hline
pair class & count & fraction \\
\hline\hline
SameResonance & 318962 & 0.26\% \\
SameWeakParent & 72371 & 0.06\% \\
SharedParton & 22638651 & 18.35\% \\
Unrelated & 100350086 & 81.33\% \\
\hline
\end{tabular}
\caption{Same-event pion pairs by ancestry.}
\label{tab:pair}
\end{table}

\begin{figure}
\includegraphics[scale=0.75,height=2in]{}
\caption{Two-particle #Delta#phi correlation decomposed by pair ancestry. Ancestry-related pairs peak near zero; unrelated pairs are flat.}
\label{}
\end{figure}
\begin{figure}
\includegraphics[scale=0.75,height=2in]{}
\caption{Two-particle #Delta#eta correlation decomposed by pair ancestry.}
\label{}
\end{figure}



\section{Mechanism ladder}


The same analysis is repeated as Pythia mechanisms are switched on one at a time (shower, MPI, colour reconnection, rope hadronization), showing how each mechanism shifts the provenance breakdown.


\begin{table}[htbp]
\centering
\begin{tabular}{l l r r r r}
\hline
group & observable & shower & shower+MPI & shower+MPI+CR & shower+MPI+CR+rope \\
\hline\hline
origin & Primary & 31.7900 & 31.6500 & 31.6700 & 31.5300 \\
origin & FromResonance & 57.9000 & 57.8800 & 57.7200 & 57.7700 \\
origin & FromWeakDecay & 10.3000 & 10.4700 & 10.6000 & 10.7000 \\
parton & LightQuark & 89.9000 & 76.2400 & 77.3600 & 79.4300 \\
parton & Strange & 6.1800 & 11.8800 & 12.6600 & 12.3500 \\
parton & Charm & 1.8500 & 2.7900 & 3.1400 & 2.9400 \\
parton & Bottom & 0.2800 & 0.2300 & 0.2600 & 0.2400 \\
parton & Gluon & 0.0000 & 8.0700 & 4.8600 & 4.3500 \\
pair & SameResonance & 0.7400 & 0.1700 & 0.2600 & 0.2300 \\
pair & SameWeakParent & 0.1600 & 0.0400 & 0.0600 & 0.0500 \\
pair & SharedParton & 39.3000 & 15.9500 & 18.3500 & 19.9100 \\
pair & Unrelated & 59.8000 & 83.8400 & 81.3300 & 79.8100 \\
\hline
\end{tabular}
\caption{Provenance breakdown across the mechanism ladder (percentages).}
\label{tab:ladder}
\end{table}

\begin{figure}
\includegraphics[scale=0.75,height=2in]{}
\caption{#Delta#phi of shared-parton pairs across the mechanism ladder.}
\label{}
\end{figure}
\begin{figure}
\includegraphics[scale=0.75,height=2in]{}
\caption{#Delta#phi of unrelated (combinatorial) pairs across the ladder.}
\label{}
\end{figure}
\begin{figure}
\includegraphics[scale=0.75,height=2in]{}
\caption{p_{T} of resonance-feed-down pions across the mechanism ladder.}
\label{}
\end{figure}
\begin{figure}
\includegraphics[scale=0.75,height=2in]{}
\caption{p_{T} of strange-ancestor pions across the mechanism ladder.}
\label{}
\end{figure}



\section{Discussion}


The resonance / weak-decay split depends on a curated resonance list; the primary fraction, taken directly from the production stage, is exact.  Sub-percent pair fractions carry statistical uncertainty that should be quantified before quotation.




\end{document}
